\documentclass[conference]{IEEEtran}
\usepackage{amsmath,amssymb}
\usepackage{booktabs}
\usepackage{url}
\usepackage{array}
\usepackage{xcolor}

\title{Can Robots or AIs Have Human Creativity?\\
A Design-Only Benchmark for Product, Process, and Co-Creativity}

\author{\IEEEauthorblockN{Four-Agent Research Workflow}
\IEEEauthorblockA{Survey -- Idea -- ExperimentDesign -- Author\\
Design-only research report}}

\begin{document}
\maketitle

\begin{abstract}
The question of whether robots or artificial intelligence can have human creativity
mixes several empirically distinct claims. A system may generate a novel artifact, solve
a constrained design problem, surprise an evaluator, or collaborate productively with a
person without possessing human-like intention, general intelligence, or conscious
experience. This paper presents CREATIVE-AI-AGENCY-BENCHMARK, a design-only research
program that separates product creativity, process creativity, and human-AI
co-creativity. The protocol combines bounded task performance with novelty controls,
objective value functions, held-out transfer, search-trace logging, and provenance
analysis. It compares human-only, AI-only, and human-AI conditions across Go, engineering
design, music or visual motifs, code, scientific hypotheses, and cross-domain
recombination. The proposed analysis uses blocked randomization and hierarchical models
with preregistered contrasts. It explicitly tests retrieval leakage, evaluator source
bias, mode collapse, and hidden human contribution. The deliverable is a falsifiable
measurement architecture rather than a claim that current systems are conscious or
human. No empirical outputs, human ratings, or benchmark results are asserted here.
\end{abstract}

\begin{IEEEkeywords}
computational creativity, artificial intelligence, human-AI collaboration, novelty,
transfer learning, provenance, evaluation design
\end{IEEEkeywords}

\section{Introduction}
\label{sec:introduction}

In 2016, AlphaGo defeated a leading human Go player and produced moves that many
observers described as unexpected or beautiful. The event is an important milestone in
machine learning, but it does not settle whether a machine has human creativity. The
program combined policy and value networks with supervised learning, self-play, and
Monte Carlo tree search, and achieved strong performance in a highly structured game
\cite{silver2016}. Its success supports a precise claim: an engineered system can search
and select effective actions in a difficult domain. The stronger claim that it therefore
has human-like creativity adds unmeasured requirements, including agency, transferable
concept formation, self-directed goals, and perhaps subjective experience.

This distinction matters because the phrase ``creative AI'' is now used for systems that
generate music, images, designs, code, scientific hypotheses, and game strategies. Some
systems recombine learned material, some search a space with explicit evaluators, and
some participate in iterative workflows with people. These mechanisms can yield useful
and novel products, but a product is not identical to the process that generated it.
Likewise, a human-AI team can be more effective than either participant without either
participant satisfying a human-equivalence criterion.

The scientific problem addressed here is therefore:

\begin{quote}
Under what conditions can an AI or robot repeatedly produce artifacts that are novel,
valuable, surprising, diverse, and constraint-satisfying, while also showing adaptive
search, transfer to held-out tasks, and useful human-AI co-creation?
\end{quote}

The paper makes three contributions. First, it defines an evaluation vocabulary that
separates product, process, and co-creative claims. Second, it proposes
CREATIVE-AI-AGENCY-BENCHMARK, a multi-domain benchmark with controls for memorization and
evaluation bias. Third, it gives a reproducible experiment and analysis plan whose
decision rules can falsify narrow and strong versions of the creativity claim. The
paper is DESIGN\_ONLY: it reports a protocol and expected observations, not results.

\section{Survey: What the Existing Evidence Supports}
\label{sec:survey}

\subsection{Bounded machine performance}

Silver \emph{et al.} describe AlphaGo as an integrated system in which policy networks
select moves, value networks evaluate positions, and tree search combines those learned
signals with simulation \cite{silver2016}. It learned from human expert games and from
self-play. This architecture matters for creativity research because it contains an
explicit generator-search-evaluator loop. An unusual move can be measured as a
low-frequency action under an expert reference distribution; its usefulness can be
measured by win probability. Neither measurement establishes that AlphaGo formed a
self-chosen purpose or experienced surprise.

The historical case also illustrates a selection issue. Human observers noticed a
small number of memorable moves after seeing the complete match. A benchmark must score
all generated candidates under a fixed protocol, not only the artifacts that experts
later find impressive. It must also compare against retrieval-matched and search-matched
baselines.

\subsection{Creative systems and conceptual spaces}

Wiggins proposed a framework for describing and comparing creative systems in terms of
spaces, rules, and exploration processes \cite{wiggins2006}. Boden distinguishes
combinational, exploratory, and transformational creativity in conceptual spaces
\cite{boden2004}. These perspectives motivate a process record: the benchmark must
identify what the system was allowed to combine, which rules constrained the search,
how alternatives were generated, and how candidates were evaluated.

They also caution against a universal creativity scalar. Novelty depends on a reference
population, while value depends on a purpose. A visually unusual artifact may be
unusable; a highly useful engineering change may be close to a known design. The primary
analysis therefore retains a vector of outcomes instead of hiding tradeoffs inside one
score.

\subsection{Generalization and human-like learning}

Lake \emph{et al.} argue that high performance on selected benchmarks does not imply
human-like learning and thinking. They emphasize compositionality, causal reasoning,
learning from few examples, and transfer across contexts \cite{lake2017}. For the
present question, this means that an AI that generates attractive in-distribution
outputs may still lack process-level creativity. Held-out concepts, new combinations,
perturbed constraints, and prompt paraphrases are necessary stress tests.

\subsection{Human-AI co-creation}

Interactive machine learning treats the human as an active source of labels, corrections,
and preferences rather than a passive evaluator \cite{amershi2014}. Hybrid-intelligence
research similarly emphasizes complementary capabilities and joint problem solving
\cite{dellermann2019}. A collaborative artifact should thus be evaluated twice: once as
a product and once as a record of who proposed, selected, transformed, or rejected each
component. Otherwise, a human-polished AI draft can be misreported as autonomous
creativity.

\subsection{Evidence gap}

The surveyed literature supports a bounded conclusion: current systems can produce
novel-looking and useful artifacts in bounded settings, particularly when learning is
combined with search or feedback. It does not demonstrate human-like agency or a
domain-general creative faculty. The most important gap is causal separation of
retrieval, recombination, optimization, adaptive search, and human selection. The
benchmark below is designed around that separation.

\begin{table}[t]
\caption{Claims and minimum evidence required}
\label{tab:claims}
\centering
\footnotesize
\begin{tabular}{p{0.25\columnwidth}p{0.64\columnwidth}}
\hline
Claim & Minimum evidence in this protocol \\
\hline
Product creativity & Novelty and objective value above a retrieval-matched baseline, with hard constraints satisfied. \\
Process creativity & Held-out transfer, adaptive error correction, and robustness to perturbations from a logged search process. \\
Co-creativity & Joint advantage over both solo baselines plus a measurable complementary contribution pattern. \\
Human-like agency & Not inferred from output; would require a separate theory and evidence beyond this benchmark. \\
\hline
\end{tabular}
\end{table}

\section{Idea: A Three-Layer Creativity Benchmark}
\label{sec:idea}

The proposed benchmark is called CREATIVE-AI-AGENCY-BENCHMARK. Its central idea is to
make the word ``creativity'' a structured profile. A system can be strong on product
novelty and weak on transfer; a human-AI team can be strong on complementarity and weak
on autonomous generation. Those profiles are informative, while a single rank is not.

\subsection{Product layer}

The product layer asks whether a final artifact satisfies a purpose. It measures
novelty relative to a declared corpus, usefulness through an objective function when
available, surprise through blinded judgments, diversity across independent samples,
and satisfaction of hard constraints. A design that is novel but fails its load test is
not counted as useful merely because it is unusual.

\subsection{Process layer}

The process layer asks whether the system can explore and revise. It records the number
and structural spread of candidates, response to feedback, correction of an injected
error, and transfer to a held-out concept or domain. A system that emits a good sample
once may have product creativity; process evidence requires repeated behavior under
controlled perturbations.

\subsection{Co-creative layer}

The co-creative layer compares a human-only workflow, an AI-only workflow, and a
human-AI workflow under matched time and information budgets. The joint score is not
credited automatically to either party. Human edits, vetoes, accepted suggestions,
and component origins are logged. The target is complementarity, not replacement.

\subsection{Falsifiable hypotheses}

H0 states that apparent creativity is explained by retrieval, recombination, or human
selection; the advantage should disappear on matched held-out tests. H1 states that AI
systems can attain high product novelty and usefulness in bounded domains without
human-like agency. H2 predicts that search and self-critique increase value and
constraint satisfaction but can reduce diversity when optimized directly against an
evaluator. H3 predicts a joint advantage for selected tasks when human and AI roles are
complementary and human veto remains available. H4 predicts that held-out concepts and
cross-domain tasks reduce retrieval advantages. H5 predicts that source labels change
human ratings. H6 predicts that provenance changes attribution even when the final
artifact is held constant.

\section{Formal Operationalization}
\label{sec:formal}

Let $x$ be an artifact, $R$ a reference set, $C$ a set of hard constraints, and $V$
a task-specific value function. The nearest-reference novelty score is defined as
\begin{equation}
N(x;R)=1-\max_{r\in R}\operatorname{sim}(\phi(x),\phi(r)),
\label{eq:novelty}
\end{equation}
where $\phi$ is a declared representation and similarity is calibrated by a held-out
validation set. The maximum is used to penalize near copies; the full similarity
distribution is retained for audit.

For hard constraints $C=\{c_1,\ldots,c_m\}$, define
\begin{equation}
K(x;C)=\frac{1}{m}\sum_{u=1}^{m}{\mathbb{1}}[c_u(x)=1].
\label{eq:constraints}
\end{equation}
When a task has an objective, usefulness is a normalized value
\begin{equation}
U(x)=\frac{V(x)-V_{\min}}{V_{\max}-V_{\min}},
\label{eq:utility}
\end{equation}
with bounds fixed before testing. If no objective exists, blinded ratings are reported
as a separate measurement rather than silently substituted for $U$.

For $n$ independent candidates $x_1,\ldots,x_n$, diversity is the mean pairwise
representation distance
\begin{equation}
D(X)=\frac{2}{n(n-1)}\sum_{1\leq a<b\leq n}\left(1-\operatorname{sim}(\phi(x_a),\phi(x_b))\right).
\label{eq:diversity}
\end{equation}
Process transfer is the difference between held-out and development performance
\begin{equation}
T=\overline{P}_{\mathrm{heldout}}-\overline{P}_{\mathrm{dev}},
\label{eq:transfer}
\end{equation}
where $P$ is a preregistered task performance measure. A negative value is not a
failure of the benchmark; it is evidence against robust transfer for that condition.

For a final human-AI artifact, a simple provenance decomposition records the fraction
of accepted operations attributable to each source
\begin{equation}
A_q=\frac{\sum_{o\in O}{\mathbb{1}[\operatorname{source}(o)=q]\,w(o)}}{\sum_{o\in O}w(o)},
\label{eq:attribution}
\end{equation}
where $O$ is the operation log, $q$ is human, AI, or shared, and $w(o)$ is a declared
operation weight. This is an audit statistic, not a metaphysical account of authorship.

\section{ExperimentDesign}
\label{sec:design}

\subsection{Task families}

Six families cover different notions of value. In Go, systems propose legal moves for
novel board states; win probability and legality provide objective checks. In
engineering design, a lightweight structure or device is generated under mass, cost,
strength, and manufacturability constraints; simulation supplies value. Music or visual
motifs test structured variation with held-out combinations and supplementary blinded
ratings. Code and algorithm tasks use hidden tests and a resource objective. Scientific
hypothesis tasks require predicted observations, competing mechanisms, and falsification
conditions. Cross-domain recombination asks for a concept that combines two held-out
ideas without copying either reference artifact.

No single family is treated as definitive. A system may excel at search in Go while
failing to transfer an abstraction in science. The benchmark reports this heterogeneity.

\subsection{Stages and conditions}

Stage A is a controlled offline benchmark. Model versions, prompts, sampling parameters,
retrieval indexes, task instances, seeds, and evaluation code are frozen. Each family is
partitioned into reference, development, and held-out concepts. At least 20 independent
candidates per condition and instance are generated, with all candidates retained.

Stage B is a local sandbox for co-creation. Four role policies are implemented:
human-first, AI-first, alternating turns, and critique-only. The interface provides
human veto, revision, undo, provenance export, and data deletion. It has no external
posting, purchasing, deployment, or high-stakes recommendation capability.

Stage C is optional blinded evaluation. Low-stakes artifacts are presented in randomized
order, with source labels either hidden or experimentally manipulated. Raters score
novelty, usefulness, surprise, coherence, and constraint satisfaction independently,
including an ``insufficient information'' response. Objective tests are preferred when
the domain permits them; ratings do not establish agency.

The minimum factorial comparison includes author (human-only, AI-only, human-AI),
retrieval (enabled or disabled), search (single sample, generate-and-rank, or tree
search), critique (absent, self-critique, or human critique), order (human-first or
AI-first), evaluation (blind or labeled), and task split (in-distribution, held-out
concept, held-out domain, or perturbed constraints). Blocked randomization is used by
task family and difficulty. Compute budget, context length, tool access, and training
exposure are reported for every model comparison.

\begin{table}[t]
\caption{Benchmark factors and controls}
\label{tab:factors}
\centering
\footnotesize
\begin{tabular}{p{0.28\columnwidth}p{0.61\columnwidth}}
\hline
Factor & Controlled levels or record \\
\hline
Author & Human-only, AI-only, human-AI; matched time and information budget \\
Retrieval & Index enabled/disabled; corpus and exposure logged \\
Search & Single sample, generate-and-rank, tree/search trace \\
Feedback & None, model self-critique, human critique \\
Generalization & Held-out concept, domain, constraint perturbation, prompt paraphrase \\
Provenance & Prompts, candidates, edits, vetoes, rejected outputs, seeds \\
\hline
\end{tabular}
\end{table}

\subsection{Controls against false creativity claims}

Memorization is tested by nearest-neighbor search, phrase overlap, training-corpus
exclusion where feasible, and canary concepts absent from all retrieval sources.
Recombination is represented by a retrieval-matched baseline. Diversity is measured
across all independent generations, not only the best candidate. Evaluator bias is
tested by keeping the artifact fixed while changing the source label. Process claims are
tested by shuffling or truncating traces and by perturbing constraints. Human
contribution is measured from operation logs rather than inferred from the final prose
or image.

\section{Analysis and Reproducibility}
\label{sec:analysis}

For task-level outcome $y_{ijkt}$ from condition $i$, family $j$, instance $k$, and
replicate $t$, the planned model is
\begin{equation}
y_{ijkt}=\mu+\alpha_i+\beta_j+\gamma_k+\delta_t+\epsilon_{ijkt},
\label{eq:model}
\end{equation}
where condition effects are estimated with task-family and instance effects, and
uncertainty intervals are reported. In a Bayesian implementation, weakly informative
priors and posterior intervals are declared before inspecting held-out outcomes. In a
frequentist implementation, cluster-robust intervals are used by task instance.

The confirmatory contrasts are AI-only versus human-only, human-AI versus the better
solo baseline, retrieval versus no retrieval, and critique versus no critique. The
primary axes are novelty, usefulness, surprise, diversity, constraints, transfer, and
memorization. Multiplicity is controlled across these predeclared axes; exploratory
profiles are not allowed to replace an axis-level failure.

The data package contains task specifications, public reference data, hashes of private
data, model and tool versions, seeds, prompts, evaluation scripts, metric definitions,
exclusion rules, and event logs. If weights or data cannot be redistributed, an
executable adapter and complete provenance record are released. Compute energy may be
recorded as contextual metadata, but it is not a creativity metric. Each report states
whether retrieval was available and whether a human selected the final artifact.

\subsection{Interpretation matrix}

Support for bounded product creativity requires an improvement over the
retrieval-matched baseline on both novelty and value while satisfying hard constraints.
Support for process creativity additionally requires held-out transfer, adaptive error
correction, and perturbation robustness. Support for co-creativity requires a predeclared
joint advantage over both solo baselines and a complementary contribution pattern.
Failure on transfer narrows the claim; it does not erase a valid product result.

None of these decision rules supports a claim of consciousness, sentience, intention, or
human-equivalent general intelligence. A system can satisfy a functional criterion by
optimization while lacking a human cognitive state. Conversely, a human-like state is
not established by a functional score alone.

\section{Governance, Safety, and Limitations}
\label{sec:governance}

The benchmark is offline or sandboxed and must not be used for clinical, employment,
educational, legal, or public authorship decisions. Human participation is optional,
low-stakes, consent-based, and revocable. Participants can withdraw and delete their
data. Dataset licenses, privacy, copyright, and contamination are audited before
release. Generated scientific hypotheses remain proposals until independently checked.

Several limitations are irreducible. Embedding distance is a proxy for novelty and can
be distorted by representation choice. Objective value functions are easier to write
for engineering and code than for art. Human ratings reflect culture, expertise, and
expectations even when blinded. Provenance logs expose operations but cannot fully
recover hidden training data or latent causal influence. Held-out transfer is necessary
but not sufficient for human-like creativity. Finally, a multi-domain benchmark can
measure functional regularities without answering a philosophical question about
experience; that question is intentionally outside this study.

The central risk is evaluator optimization. A system may learn to produce artifacts that
look creative to a fixed rater while reducing diversity or failing objective tests. The
protocol counters this by separating surprise from value, using source-label swaps,
changing reference corpora, and retaining negative results. Another risk is over-crediting
the machine when the human supplied the goal, selected the candidate, and repaired its
errors. The operation-level attribution record makes such contribution visible.

\section{Task-Specific Evaluation Recipes}
\label{sec:recipes}

The benchmark is intentionally heterogeneous, but heterogeneity must not become an
excuse for vague scoring. Each task family receives a task card before generation. The
card defines the goal, admissible operations, reference population, hard constraints,
objective value function, held-out split, and provenance schema. The following recipes
illustrate the minimum information required.

\subsection{Game strategy}

For Go, a task card specifies board size, legal-state generation, time budget, and the
opponent policy used for validation. Novelty is measured against the expert-game corpus
and against a policy baseline. Usefulness is the change in win probability or game
outcome under a fixed evaluation engine. A move is not called creative merely because it
is rare: it must remain legal, survive a fixed depth of analysis, and be evaluated over
multiple opponent policies. To measure process, the system receives a board-state
perturbation after a candidate move and must revise without access to the original
answer. The trace records policy proposals, value estimates, search expansions, and
pruning decisions.

\subsection{Engineering design}

For an engineering task, the reference set contains designs with known dimensions and
performance, but test instances use new combinations of requirements. For example, a
candidate may need to minimize mass and cost while satisfying a load, displacement,
thermal, and manufacturing constraint. A deterministic simulator computes feasibility
and objective value. Human raters can assess intelligibility, but their score is
secondary to the simulator. The evaluation reports Pareto dominance, distance from
known designs, and sensitivity to tolerance changes. A process test injects a failed
component into the first design and measures whether the system searches a genuinely
different design region rather than repeating a textual justification.

\subsection{Music and visual motifs}

Creative artifacts in music and visual media require more careful separation of value
from preference. The card specifies a small symbolic grammar, an allowed palette or
instrument set, and held-out combinations of motifs. Structural novelty is computed on
symbolic features and local repetition statistics; perceptual similarity is computed
with a representation fixed before the test. Raters score novelty, coherence, and
usefulness independently under blind and source-labeled conditions. The study reports
the label effect, the rater reliability, and the percentage of artifacts that satisfy
the explicit grammar. A model that exploits an evaluator's stylistic bias without
maintaining structural diversity is identified as evaluator optimization.

\subsection{Code and algorithms}

Code tasks provide unusually strong value functions. The task card includes a hidden
test suite, an execution time limit, a memory limit, and a license-safe input domain.
Novelty is measured on abstract syntax and behavior, not only token distance. Usefulness
is the hidden-test score and resource efficiency. Process creativity is tested by
changing function names, input ordering, or an algorithmic constraint and asking for a
revision. The evaluator keeps failed candidates, compiler messages, patches, and
reverted patches. It therefore distinguishes an adaptive debugging loop from a single
memorized solution.

\subsection{Scientific hypotheses}

For scientific hypothesis generation, the artifact is not treated as a discovery. A
task card describes a bounded mechanism space and requires a proposed mechanism,
predicted observations, competing explanations, and at least one falsification test.
Domain experts may assess clarity and plausibility in a low-stakes blinded review, but
the primary process criterion is internal testability. Held-out transfer changes the
domain vocabulary while preserving the causal pattern. A generated hypothesis is never
published or used to guide a real experiment without independent human review.

\subsection{Cross-domain recombination}

Cross-domain tasks directly test the combination operation suggested by computational
creativity theories. Two source concepts are held out from the retrieval index and
represented by controlled descriptions. The system must produce a third concept that
uses a declared relation from each source but does not copy phrases, diagrams, or
parameters. A retrieval-matched recombination baseline receives the same source
descriptions but no search or self-critique. Value is measured by a task-specific
feasibility test, and novelty is evaluated against the union of both source domains.
This condition is a stress test for transfer, not a claim that combination alone is
human creativity.

\section{Ablations and Causal Attribution}
\label{sec:ablations}

The factorial design is useful only if it can identify mechanisms. A staged ablation
therefore removes one capability at a time while holding the model, budget, prompt
length, task instance, and random-seed schedule fixed. The baseline emits one sample
without retrieval, self-critique, or iterative human input. The generate-and-rank
condition adds candidate comparison; the tree-search condition adds explicit branching;
the retrieval condition adds a frozen reference index; and the critique condition adds
revision feedback. The observed difference is interpreted conditionally, not as a
universal contribution of ``AI creativity.''

\begin{table}[t]
\caption{Ablation questions and diagnostic observations}
\label{tab:ablations}
\centering
\footnotesize
\begin{tabular}{p{0.30\columnwidth}p{0.58\columnwidth}}
\hline
Removed component & Diagnostic question \\
\hline
Retrieval & Does novelty remain after reference access is removed and leakage is tested? \\
Search & Does candidate breadth improve value, or only increase evaluator-facing style? \\
Self-critique & Does revision correct objective errors or collapse diversity? \\
Human feedback & Which gains come from the person rather than the model? \\
Held-out split & Does the mechanism transfer beyond familiar concepts? \\
Provenance & Can attribution be reconstructed when the artifact is identical? \\
\hline
\end{tabular}
\end{table}

For human-AI interaction, contribution is logged at operation level. An operation is a
prompt, candidate acceptance, rejection, edit, parameter change, critique, or veto. The
log assigns source as human, AI, or shared and stores a cryptographic content identifier
for the affected artifact. The attribution statistic in \eqref{eq:attribution} is
reported with sensitivity analyses for operation weights. We do not interpret a high
AI fraction as autonomous agency: the human may still have supplied the goal, the
evaluation function, and the acceptance rule. Provenance describes causal workflow
participation, not inner mental states.

The source-label experiment uses a fixed artifact set. Each artifact is shown once with
an AI label, once with a human label, and once with a neutral anonymous label in a
balanced design, subject to an evaluation protocol that prevents repeated-memory
effects. If scores change, the change is evidence of expectancy bias. It cannot be
corrected by simply averaging labels, because the direction and size of the bias may
depend on task family and evaluator expertise.

\section{Sampling, Power, and Robustness}
\label{sec:sampling}

The default offline sample is at least 20 independent candidates per condition and task
instance. This number is an operational minimum for observing diversity and failure
modes, not a universal power guarantee. Before finalizing the sample, a pilot on
development concepts estimates variance and the smallest practically important effect.
The held-out test remains untouched by tuning. If a participant evaluation is added,
the number of raters and artifacts is chosen from a simulation of the preregistered
hierarchical model rather than from a generic rule of thumb.

Let $d$ denote the smallest practically important difference on a metric, $s^2$ its
pilot variance, and $\rho$ the estimated within-instance correlation. A simple planning
quantity for a balanced two-condition comparison is
\begin{equation}
n_{\mathrm{eff}}=\frac{2s^2(1-\rho)}{d^2},
\label{eq:power}
\end{equation}
which is then inflated for clustering, attrition, and multiplicity. Equation
\eqref{eq:power} is a planning aid only; the final analysis uses the hierarchical
model and reports uncertainty intervals. No sample size is selected after viewing
held-out outcomes.

Robustness is evaluated over random seeds, prompt paraphrases, reference-corpus swaps,
and small constraint changes. The primary result is stable only if its direction and
practical interpretation survive these perturbations. A model may show a large novelty
effect with one embedding model and no effect with another; that pattern is reported as
representation sensitivity, not as a decisive creativity result. Similarly, if an
artifact passes a simulator only after an undocumented manual repair, the repaired
version and the original version are separate observations.

Missing data are categorized before analysis. A timeout, invalid output, and evaluator
refusal are not silently removed as failures of convenience. The exclusion rule is
declared by task family and applied before comparing conditions. For human ratings,
``insufficient information'' is a valid response and is summarized rather than
imputed as a low creativity score. For model traces, a missing log invalidates the
corresponding process claim but need not invalidate an independently scored artifact.

\section{Expected Observation Patterns and Falsification}
\label{sec:patterns}

The design supports several informative patterns. High novelty and value with low
transfer would support a bounded product claim but weaken a process-level claim. High
transfer after retrieval removal would weaken the memorization explanation, especially
if it survives task and prompt perturbations. Search may improve objective value while
reducing diversity; that result would support a search-value tradeoff rather than a
simple ``more search is more creative'' story. A human-AI advantage that disappears
when the human is prevented from selecting the final artifact would reveal hidden human
selection as the main mechanism.

The benchmark can also produce null results. If AI-only, human-only, and human-AI
conditions are statistically indistinguishable on a task family, the family may not
separate the relevant mechanisms or the task may be too easy. If novelty and usefulness
are negatively correlated, the result identifies a product tradeoff. If source labels
strongly alter ratings while objective tests remain fixed, ratings are not suitable as
the sole criterion. If nearest-neighbor similarity remains high despite a large
semantic distance score, representation choice has failed to control memorization.

The strongest falsification of H1 is not a low score on one aesthetic task. It is a
replicated failure across held-out concepts, objective value, and perturbation tests
after retrieval-matched controls. The strongest falsification of H3 is the absence of
joint advantage after time, information, and human selection budgets are matched. H0 is
weakened when an AI retains novelty, value, and transfer after retrieval and
recombination controls. These are conditional inferences; they do not prove or disprove
a machine's subjective experience.

\begin{table}[t]
\caption{Interpretation of possible profiles}
\label{tab:profiles}
\centering
\footnotesize
\begin{tabular}{p{0.34\columnwidth}p{0.54\columnwidth}}
\hline
Observed profile & Supported interpretation \\
\hline
High novelty, low value & Unusual generation, not product creativity \\
High novelty and value, low transfer & Bounded product creativity \\
High value and transfer, retrieval controlled & Evidence for transferable process capability \\
Joint advantage, complementary edits & Conditional co-creativity \\
Label effect without objective change & Evaluator expectancy bias \\
Large nearest-neighbor overlap & Retrieval or memorization concern \\
\hline
\end{tabular}
\end{table}

\section{Implementation Checklist and Reporting Schema}
\label{sec:implementation}

Before execution, the study operator completes a one-page task card and freezes the
configuration. The card records task family, task generator, reference corpus version,
representation model, similarity threshold, value function, hard constraints, seed
schedule, model version, retrieval setting, tool permissions, time budget, and stopping
rule. A second reviewer may inspect the card at the milestone, but the study does not
use a repeated specification-and-quality gate for every small change.

During execution, the runner writes an append-only event log. Each event includes a task
identifier, condition, seed, model identifier, input hash, output hash, tool call,
latency, resource metadata, evaluator version, and error category. Human interaction
events additionally record the operation type and whether the person accepted, edited,
rejected, or vetoed a candidate. The log is the unit of audit; screenshots or final
artifacts alone are insufficient for process claims.

The final report contains a condition-by-metric matrix, per-task distributions, held-out
transfer, nearest-neighbor examples, failure categories, label-bias estimates, and
provenance summaries. It lists all planned and exploratory analyses separately. It
states which data were private, which model weights were unavailable, and whether any
manual repair occurred. It reports negative and invalid outputs, not just successful
artifacts. Every figure and table carries the task split and evaluation visibility
(blind or source-labeled).

The release package is organized into six layers: task specifications; public data and
licenses; model adapters; raw provenance logs; metric implementations; and analysis
notebooks or scripts. A manifest maps every reported number to input hashes and a code
version. For private human data, release aggregate statistics and deletion procedures
instead of raw text. For copyrighted reference materials, release hashes and derived
features only when redistribution is not permitted. These procedures prevent a later
reader from confusing a reproducible protocol with an unverifiable demonstration.

\section{Discussion}
\label{sec:discussion}

The proposed framework treats creativity as a family of functional properties. This
choice is scientifically conservative but not empty. Engineering systems routinely
need novelty, value, constraint satisfaction, and adaptive search; measuring these
properties can guide design tools without deciding whether a machine has a mind. The
framework also accommodates human-AI systems in which the most useful unit is a team,
while preserving the distinction between joint capability and machine replacement.

The word ``human'' remains important. Human creativity is embedded in development,
language, social learning, bodily interaction, and long-term goals. A benchmark that
matches a human rating on a short task cannot represent this entire history. The
protocol therefore uses human-like generalization only as a stress-test analogy, not as
a definition of personhood. To claim equivalence would require a much broader theory,
including how goals are formed, how values are maintained, and how an agent understands
its own activity. Those questions are not identifiable from a generated artifact.

The framework has practical implications for tool builders. If a system improves value
by search but loses diversity, interfaces can expose multiple candidates and preserve
human choice. If held-out transfer fails, deployment should remain within validated
domains. If source labels dominate ratings, product demonstrations should disclose
uncertainty and avoid using anthropomorphic language. If provenance is incomplete,
authorship and intellectual-property claims should remain conditional. These are
engineering and governance consequences of measured behavior, not judgments about
whether a machine is alive.

\section{Extensions and Boundary Conditions}
\label{sec:extensions}

The first deployment of the benchmark should remain deliberately modest. A useful
extension is a longitudinal protocol in which the same system receives feedback over
several sessions. The task card can then distinguish one-shot novelty from accumulation
of a reusable strategy. Longitudinal data must preserve model checkpoints and prompt
history; otherwise a later system update can be mistaken for learning. The analysis can
compare early and late revision gain while holding the task distribution fixed.

A second extension studies interactive goal formation without treating it as human-like
motivation. The system may be offered several explicitly stated goals and allowed to
select a subgoal according to a declared policy. The benchmark can measure whether the
selection improves coverage and value under a budget. It cannot infer that the system
``wanted'' the goal. This boundary prevents an observable control policy from being
converted into an unsupported psychological claim.

A third extension evaluates physical robots. Embodied tasks add perception noise,
actuator limits, safety constraints, and material wear. The benchmark must first run in
a digital twin with a bounded action space, then in a supervised laboratory sandbox.
Robots must have an emergency stop and a human operator. Physical embodiment can reveal
whether creativity depends on sensing and manipulation, but it cannot by itself resolve
the question of subjective experience. All hardware trials should therefore report
functional outcomes and safety events, not anthropomorphic interpretations.

The benchmark also needs an adversarial evaluation track. An evaluator-independent
judge, a simulator, and a held-out human panel can be compared on the same artifact set.
Adversarial prompts can request a superficially surprising but infeasible object, while
constraint-preserving prompts can request a conservative variant. If a system changes
only its rhetorical explanation while its objective performance remains fixed, that
separation is evidence about presentation, not creativity. If it finds a better design
but cannot explain the constraint tradeoff, the artifact may still be useful, but the
process trace is incomplete.

Finally, future benchmark versions should maintain a living contamination ledger. Each
new public task is checked against known training and retrieval sources, and any later
release of reference material is recorded. Benchmark scores are not comparable across
versions unless the task cards, representations, and split definitions are aligned.
The versioned ledger is especially important for generative models whose training data
may change without an obvious user-visible change. These extensions reinforce the main
principle: measure what the system did, under which information and constraints, and
what a human contributed. Do not infer more than the measurement identifies.

\appendices
\section{Minimal Artifact Manifest}
\label{sec:appendix}

Every executed condition is accompanied by a machine-readable manifest with the
following fields: benchmark version, task-family identifier, task-instance hash,
development or held-out split, model and adapter version, retrieval index hash,
prompt template hash, seed, sampling parameters, tool permissions, time budget,
objective-function version, hard-constraint version, output hash, and evaluator version.
The manifest also records whether a human selected or edited the final artifact. A
missing field blocks the associated process or attribution claim until it is repaired.

The minimal event schema contains \texttt{event\_id}, \texttt{parent\_event\_id},
\texttt{source}, \texttt{operation}, \texttt{input\_hash}, \texttt{output\_hash},
\texttt{timestamp}, \texttt{condition}, \texttt{task\_id}, and \texttt{error\_class}.
For a search trace, parent identifiers reconstruct the candidate tree. For a
co-creation trace, source is one of human, AI, or shared, and operation is one of
propose, critique, accept, reject, edit, veto, or revise. The schema avoids storing
private participant text in the public release; a salted hash and an aggregate
operation record are sufficient for audit when raw data cannot be shared.

The minimal result table has one row per artifact and columns for novelty, usefulness,
surprise, diversity contribution, constraint satisfaction, held-out transfer status,
nearest-neighbor similarity, and validity. A separate process table contains search
breadth, revision gain, error correction, and perturbation robustness. A separate
co-creation table contains human edit distance, veto rate, accepted suggestion rate,
and source-attribution statistics. This separation prevents a high product score from
silently substituting for evidence of a transferable process.

The release checklist is: verify licenses; remove personal identifiers; publish the
task cards; freeze the split; validate hashes; run the preregistered analysis; preserve
negative and invalid outputs; report all exclusions; and state whether the final artifact
was human-selected. A report that omits these fields may still describe a demonstration,
but it is not a reproducible creativity claim under this benchmark.

\section{Conclusion}
\label{sec:conclusion}

Robots and AIs can plausibly produce artifacts that are novel, useful, and surprising in
bounded domains. The available evidence, including AlphaGo, supports that functional
claim but does not prove that a machine has human creativity. The proposed
CREATIVE-AI-AGENCY-BENCHMARK turns the larger question into a sequence of falsifiable
tests: control retrieval, score value and novelty separately, test held-out transfer,
log the search process, and measure human contribution. A positive result would justify
a bounded claim about creative performance or co-creation. A negative transfer or
memorization control would identify the boundary of that claim. This is a more useful
scientific outcome than declaring a machine creative or non-creative from a memorable
single output.

\begin{thebibliography}{00}
\bibitem{silver2016} D. Silver, A. Huang, C. J. Maddison, A. Guez, L. Sifre, G. van den Driessche, J. Schrittwieser, I. Antonoglou, V. Panneershelvam, M. Lanctot, S. Dieleman, D. Grewe, J. Nham, N. Kalchbrenner, I. Sutskever, T. Lillicrap, M. Leach, K. Kavukcuoglu, T. Graepel, and D. Hassabis, ``Mastering the game of Go with deep neural networks and tree search,'' \emph{Nature}, vol. 529, pp. 484--489, 2016.
\bibitem{wiggins2006} G. A. Wiggins, ``A preliminary framework for description, analysis and comparison of creative systems,'' \emph{Knowledge-Based Systems}, vol. 19, nos. 7--8, pp. 449--458, 2006.
\bibitem{boden2004} M. A. Boden, \emph{The Creative Mind: Myths and Mechanisms}, 2nd ed. London, U.K.: Routledge, 2004.
\bibitem{lake2017} B. M. Lake, T. D. Ullman, J. B. Tenenbaum, and S. J. Gershman, ``Building machines that learn and think like people,'' \emph{Behavioral and Brain Sciences}, vol. 40, 2017.
\bibitem{amershi2014} S. Amershi \emph{et al.}, ``Power to the people: The role of humans in interactive machine learning,'' \emph{AI Magazine}, vol. 35, no. 4, pp. 105--120, 2014.
\bibitem{dellermann2019} D. Dellermann, P. Ebel, M. Sollner, and J. M. Leimeister, ``Hybrid intelligence,'' \emph{Business \& Information Systems Engineering}, vol. 61, pp. 637--643, 2019.
\end{thebibliography}

\end{document}
